\input{../../../templates/es/preambulo.tex}

% Los listados se toman de src/ con \lstinputlisting; la diapositiva es el archivo.
\lstset{
    basicstyle=\ttfamily\fontsize{8pt}{9.6pt}\selectfont,
    breaklines=true,
    breakatwhitespace=true,
    columns=fullflexible,
    keepspaces=true,
    xleftmargin=0.3em,
    framesep=2pt,
    aboveskip=0pt,
    belowskip=0pt
}

\newcommand{\marcoresultado}[3]{%
    \begin{frame}{#1}
        \begin{center}
            \includegraphics[height=0.62\textheight,width=\textwidth,keepaspectratio]{#2}
        \end{center}
        \vspace{0.25em}
        {\small #3\par}
    \end{frame}%
}

\newcommand{\marcoafirmacion}[2]{%
    \begin{frame}{#1}
        {\small #2\par}
    \end{frame}%
}

\date{}

\begin{document}

\tituloleccion{07}{Ajuste de Parámetros}

\section{Esta lección}

\begin{frame}{Esta lección}
La Lección 06 construyó el instrumento de tasa de éxito y el contador de evaluaciones. Hoy esas herramientas se apuntan a los parámetros escritos arriba de cada archivo hasta ahora.

\vspace{0.6em}
\begin{itemize}
    \item Afinar desde corridas sueltas es leer ruido.
    \item La moneda honesta es el margen sobre la búsqueda ciega a igual presupuesto.
    \item Un ajuste fijo es un compromiso --- por eso existe la Lección 12.
\end{itemize}
\end{frame}

% ================= afinacion_01_corridas_sueltas =================
\begin{frame}{Evidencia emparejada}
El mismo índice de semilla se usa en ambas configuraciones:
\[
d_i=I_{i,B}-I_{i,A}\in\{-1,0,1\},\qquad
\widehat\Delta=\bar d,\qquad
\widehat{\operatorname{EE}}(\widehat\Delta)=\frac{s_d}{\sqrt n}.
\]
Este error emparejado sustituye la aproximación independiente. El mapa de 36
celdas es exploratorio para este paisaje y semillas, no un óptimo universal ni
un estudio confirmatorio de comparaciones múltiples.
\end{frame}

\section{El metodo del libro, fielmente}

\begin{frame}{El metodo del libro, fielmente}
El Capitulo 7 es el capitulo de laboratorio del libro - 42 figuras, tres
perillas (probabilidad de cruza, probabilidad de mutacion, tamano de poblacion),
cada una girada en su propio paisaje, cada configuracion mostrada como una
corrida, generacion por generacion. Es la forma natural de afinar, y es
exactamente el metodo que la Leccion 06 demolio. Este paso mantiene el
protocolo del libro y sus numeros - poblacion 16, probabilidad de mutacion 0.2,
probabilidades de cruza 0.0, 0.4, 0.7, semilla 63 - en el paisaje del curso, y
lo lee de la forma en que lo lee el libro. Luego lo lee siete veces mas.
\end{frame}

\marcoresultado{El metodo del libro, fielmente}{../figuras/afinacion_01_corridas_sueltas.png}{%
    Una corrida por config corona pc = 0.7 por 0.0021; sobre 8 semillas las victorias van 0/3/5, la división más estrecha es 0.000004 (semilla 3) y la más ancha 0.6284 (semilla 1) — el ranking es inestable}

% ================= afinacion_02_el_instrumento =================
\section{El instrumento}

\begin{frame}{El instrumento}
El Paso 1 corrio el metodo del libro y lo vio fallar: una corrida por configuracion, un
ranking decidido en el cuarto decimal, conteos de victorias que cambian con la semilla.
La Leccion 06 construyo la cura - una tasa de exito sobre cientos de corridas, juzgada desde
afuera por una referencia de cuadricula densa y una tolerancia establecida. Este paso apunta ese
instrumento a la primera perilla, la probabilidad de cruza, en la linea base
del curso: poblacion 10, probabilidad de mutacion 0.1, el paisaje de seno de la
Leccion 01, diez generaciones.
\end{frame}

\begin{frame}[fragile]{(1) las perillas son argumentos}
\lstinputlisting[basicstyle=\ttfamily\fontsize{6pt}{7.2pt}\selectfont,firstline=113,lastline=142]{../src/afinacion_02_el_instrumento.py}
\end{frame}

\begin{frame}[fragile]{(2) error\_estandar()}
\lstinputlisting[basicstyle=\ttfamily\fontsize{8pt}{9.6pt}\selectfont,firstline=159,lastline=167]{../src/afinacion_02_el_instrumento.py}
\end{frame}

\begin{frame}[fragile]{(3) medir()}
\lstinputlisting[basicstyle=\ttfamily\fontsize{7pt}{8.5pt}\selectfont,firstline=170,lastline=188]{../src/afinacion_02_el_instrumento.py}
\end{frame}

\begin{frame}[fragile]{(4) la linea base del curso}
\lstinputlisting[basicstyle=\ttfamily\fontsize{8pt}{9.6pt}\selectfont,firstline=46,lastline=46]{../src/afinacion_02_el_instrumento.py}
\end{frame}

\marcoresultado{El instrumento}{../figuras/afinacion_02_el_instrumento.png}{%
    La tasa observada sube 55.6\% → 83.0\%; la brecha de 27.4 puntos se reporta con el error estándar de diferencias 0/1 emparejadas. A la tabla aún le falta la columna de costo}

% ================= afinacion_03_la_perilla_del_presupuesto =================
\section{La perilla del presupuesto}

\begin{frame}{La perilla del presupuesto}
El Paso 2 midio la perilla de cruza y encontro una subida monotona limpia: 55.6\%
en 0.0 a 83.0\% en 1.0. Cerro con la columna faltante. La cruza no
viene gratis - cada par cruzado crea dos nuevos individuos, y cada
nuevo individuo es una evaluacion de aptitud, la unica moneda que gasta un AG. Asi
que el dial que "mejora" el algoritmo tambien infla su presupuesto, y la
pregunta se vuelve: ¿cuanto de la subida es mejor busqueda, y cuanto es
simplemente mas gasto?
\end{frame}

\begin{frame}[fragile]{(1) el contador de evaluaciones}
\lstinputlisting[basicstyle=\ttfamily\fontsize{7pt}{8.5pt}\selectfont,firstline=67,lastline=88]{../src/afinacion_03_la_perilla_del_presupuesto.py}
\end{frame}

\begin{frame}[fragile]{(2) exito\_ciego()}
\lstinputlisting[basicstyle=\ttfamily\fontsize{8pt}{9.6pt}\selectfont,firstline=195,lastline=205]{../src/afinacion_03_la_perilla_del_presupuesto.py}
\end{frame}

\begin{frame}[fragile]{(3) la columna de margen}
\lstinputlisting[basicstyle=\ttfamily\fontsize{8pt}{9.6pt}\selectfont,firstline=221,lastline=230]{../src/afinacion_03_la_perilla_del_presupuesto.py}
\end{frame}

\marcoresultado{La perilla del presupuesto}{../figuras/afinacion_03_la_perilla_del_presupuesto.png}{%
    Cruza 0.0 gasta 20 evals/corrida, 1.0 gasta 120; el margen sobre ciego cae +30.6 → +0.8 puntos; las 100 evaluaciones extra le compran al AG 27.4 puntos donde 100 extracciones ciegas compran 57.2 — la perilla es de presupuesto}

% ================= afinacion_04_mutacion =================
\section{Mutacion, y donde se esconde la curva de campana}

\begin{frame}{Mutacion, y donde se esconde la curva de campana}
Todos los libros de texto dibujan la curva de mutacion como una campana: muy poca mutacion y
la busqueda se estanca, demasiada y degenera en una caminata aleatoria, con un
punto optimo en el medio. El Capitulo 7 de Gridin muestra corridas sueltas en 0.0, 0.3 y
1.0 y hace gestos hacia la misma forma. Este paso mide la perilla correctamente -
siete configuraciones, 500 corridas cada una, cruza fija en 0.8 - y la curva de campana
no esta ahi. La tasa de exito bruta sube monotonicamente hasta
la mutacion 1.0.
\end{frame}

\begin{frame}[fragile]{(1) la segunda perilla VALORES\_PM barridos, cruza fija en 0.8}
\lstinputlisting[basicstyle=\ttfamily\fontsize{8pt}{9.6pt}\selectfont,firstline=199,lastline=204]{../src/afinacion_04_mutacion.py}
\end{frame}

\marcoresultado{Mutacion, y donde se esconde la curva de campana}{../figuras/afinacion_04_mutacion.png}{%
    La tasa bruta es monótona hasta mutación 1.0 (61.8\% → 89.8\%) — no hay curva de campana; el margen se dobla: −10.9\% en 0.0 (peor que dados), pico +5.9\% cerca de 0.2, −3.7\% en 1.0 — la curva vive en el margen}

% ================= afinacion_05_poblacion =================
\section{Poblacion, la perilla de presupuesto mas pura}

\begin{frame}{Poblacion, la perilla de presupuesto mas pura}
La cruza y la mutacion al menos fingen tratar sobre estrategia de busqueda.
El tamano de la poblacion no finge nada: una poblacion mas grande es mas evaluaciones,
punto. Es la perilla del presupuesto sin ningun disfraz, y el paso 4 de la Leccion 06
ya la atrapo una vez - la efectividad sube, la eficiencia cae. Este paso
la vuelve a medir dentro del estudio de afinacion, con las configuraciones del propio
libro (6, 10, 20, 50, mutacion 0.2 del libro), para cerrar el patron antes del
mapa final.
\end{frame}

\begin{frame}[fragile]{(1) la tercera perilla TAMANOS\_POBLACION barridos, cruza 0.8 y mutacion 0.2 fijas}
\lstinputlisting[basicstyle=\ttfamily\fontsize{8pt}{9.6pt}\selectfont,firstline=200,lastline=206]{../src/afinacion_05_poblacion.py}
\end{frame}

\marcoresultado{Poblacion, la perilla de presupuesto mas pura}{../figuras/afinacion_05_poblacion.png}{%
    Población 50 registra 500/500 éxitos (intervalo de Wilson al 95\%: 99.2\%--100\%) a 550 evals/corrida; el rendimiento colapsa 9.25 → 1.82 éxitos por 1000 evals; el margen hace pico en población 10 (+5.9\%) y desaparece en 50 — "¿cuál es mejor?" no tiene respuesta sin moneda}

% ================= afinacion_06_la_cuadricula =================
\section{La cuadricula, y el veredicto}

\begin{frame}{La cuadricula, y el veredicto}
Los pasos 2 a 5 giraron una perilla a la vez y cada perilla confeso ser una
perilla de presupuesto. La pregunta final honesta es si interactuan - si
la mejor cruza depende de la configuracion de mutacion, de la manera que los libros de texto
asumen cuando recomiendan pares como "0.8 y 0.1". Este paso mide la
cuadricula completa: seis configuraciones de cruza por seis configuraciones de mutacion, poblacion
10, 100 corridas por celda, cada celda valorada contra la busqueda ciega en su propio
presupuesto.
\end{frame}

\begin{frame}[fragile]{(1) la cuadricula}
\lstinputlisting[basicstyle=\ttfamily\fontsize{8pt}{9.6pt}\selectfont,firstline=192,lastline=204]{../src/afinacion_06_la_cuadricula.py}
\end{frame}

\begin{frame}[fragile]{(2) CORRIDAS = 100 cambiado de 500: treinta y seis celdas tienen que encajar en el reloj}
\lstinputlisting[basicstyle=\ttfamily\fontsize{8pt}{9.6pt}\selectfont,firstline=47,lastline=47]{../src/afinacion_06_la_cuadricula.py}
\end{frame}

\marcoresultado{La cuadricula, y el veredicto}{../figuras/afinacion_06_la_cuadricula.png}{%
    La esquina (0, 0) recupera la búsqueda ciega (9.0\% ± 2.9\% vs 13.4\% para 10 extra ciegos — el instrumento está calibrado); la mejor celda (1.0, 0.5) puntúa 90.0\% a 160 evals con margen de −0.01\%; solo 20 de 36 celdas superan su presupuesto — los operadores se ganan el sueldo en las baratas}

\section{Siguiente}

\begin{frame}{Siguiente}
La población es la perilla de presupuesto sin disfraz: una tasa de éxito perfecta se puede comprar, a costa de un colapso en el rendimiento por evaluación.

\vspace{0.8em}
\textbf{Lección 08 --- Funciones de caja negra}

Hasta ahora cada paisaje se podía dibujar y recorrer por fuerza bruta. La 08 quita la imagen. Un ajuste fijo como compromiso es la 12.
\end{frame}
\end{document}
